\documentclass[11pt,letterpaper]{article}

% --- Paquetes Esenciales de Matemática y Tipografía ---
\usepackage[utf8]{inputenc}
\usepackage[spanish,es-tabla]{babel}
\usepackage{amsmath,amssymb,amsfonts,amsthm}
\usepackage{geometry}
\geometry{top=2cm,bottom=2cm,left=2cm,right=2cm}
\usepackage{fancyhdr}
\usepackage{enumitem}
\usepackage{tasks}
\usepackage{tcolorbox}
\tcbuselibrary{skins,breakable}
\usepackage{tabularx}
\usepackage{xcolor}

% --- Paleta de Colores Académica Neutra ---
\definecolor{primaryDark}{RGB}{30, 41, 59}
\definecolor{accentBlue}{RGB}{37, 99, 235}
\definecolor{neutralGray}{RGB}{241, 245, 249}

% --- Configuración de Encabezados y Pies ---
\pagestyle{fancy}
\fancyhf{}
\renewcommand{\headrulewidth}{0.6pt}
\fancyhead[L]{\small\textbf{\textsf{DEPARTAMENTO DE MATEMÁTICA}}}
\fancyhead[R]{\small\textsf{\today}}
\fancyfoot[C]{\small\thepage}

% --- Cajas Didácticas ---
\newtcolorbox{cajainstruccion}{
  colback=neutralGray,
  colframe=primaryDark,
  arc=3mm,
  boxrule=0.8pt,
  fonttitle=\bfseries\sffamily,
  coltitle=white,
  colbacktitle=primaryDark,
  title=Instrucciones Generales
}

\newtcolorbox{cajaejemplo}[1]{
  colback=white,
  colframe=accentBlue,
  arc=2mm,
  boxrule=1pt,
  fonttitle=\bfseries\sffamily,
  title=#1
}

\begin{document}

% --- Membrete Estándar Institucional ---
\noindent
\begin{tabularx}{\textwidth}{|X|l|l|}
\hline
\textbf{Estudiante:} & \textbf{Curso:} & \textbf{Fecha:} \\
\hline
\textbf{Evaluación:} <<METADATA_TITLE>> & \textbf{Puntaje Total:} <<TOTAL_POINTS>> pts & \textbf{Nota:} \\
\hline
\end{tabularx}

\vspace{0.4cm}

\begin{cajainstruccion}
Lea atentamente cada situación problema. Justifique formalmente sus procedimientos matemáticos. En los ítems de selección múltiple, marque únicamente la opción correcta.
\end{cajainstruccion}

\vspace{0.4cm}

% --- Cuerpo Generado Dinámicamente por el Formatter ---
<<CONTENT_BODY>>

\end{document}
